\documentclass{beamer}

% You can also use a 16:9 aspect ratio:
%\documentclass[aspectratio=169]{beamer}
\usetheme{TACC16}

% It's possible to move the footer to the right:
%\usetheme[rightfooter]{TACC16}


\begin{document}
\title[XALT]{XALT: Understanding HPC Usage via Job Level Collection}
\author{Robert McLay} 
\date{June 25, 2018} 

% page 1
\frame{\titlepage} 

% page 2
\begin{frame}{XALT: What runs on the system}
  \begin{itemize}
    \item A U.S. NSF Funded project: PI: Mark Fahey and Robert McLay
    \item A Census of what programs and libraries are run
    \item Running at TACC, NICS, U. Florida, KAUST, ...
    \item Integrates with TACC-Stats.
    \item Has commercial support from Ellexus 
  \end{itemize}
\end{frame}

% page 3
\begin{frame}{History of XALT}
  \begin{itemize}
    \item Mark Fahey (was NICS, now ANL): ALT-D (MPI only)
    \item Robert McLay (TACC) Lariat (MPI only)
    \item Reuben Budiardja (was NICS now ORL)
    \item ALT-D $+$ Lariat $\Rightarrow$ XALT 1: (MPI only)
    \item XALT 2: All programs
  \end{itemize}
\end{frame}

% page 4
\begin{frame}{Design Goals}
  \begin{itemize}
    \item Be extremely light-weight
    \item Provide provenance data: How?
    \item How many use a library or application?
    \item Collect Data into a Database for analysis.
    \item Minimum Hardware Requirements $\Rightarrow$ Cluster and a VM
  \end{itemize}
\end{frame}

% page 5
\begin{frame}{Understanding what your users are doing}
  \begin{itemize}
    \item What programs, libraries are your users using?
    \item What are the top programs by node-hours? by counts?
    \item Are they building their own programs or using someone elses?
    \item Are Executables implemented in C/C++/Fortran?
    \item Track MPI: tasks? nodes?
    \item Track Threading via \$OMP\_NUMTHREADS
  \end{itemize}
\end{frame}

% page 6
\begin{frame}{Design: Linker}
  \begin{itemize}
    \item XALT wraps the linker to enable tracking of exec's
    \item The linker (ld) wrapper intercepts the user link line.
    \item Generate assembly code: key-value pairs
    \item Capture tracemap output from ld
    \item Transmit collected data in *.json format
    \item Adds codes that executes before main() and after main()
      completes
    \item Knows about SLURM, PBS, LSF, SGE
    \item Translates SLURM\_NNODES or PBS\_NUM\_NODES $\Rightarrow$ num\_nodes
  \end{itemize}
\end{frame}

% page 7
\begin{frame}{Design: Transmission to DB}
  \begin{itemize}
    \item File: collect nightly
    \item Syslog: Use Syslog filtering (or ELK)
    \item Future: RabbitMQ
  \end{itemize}
\end{frame}

% page 8
\begin{frame}{Lmod to XALT connection}
  \begin{itemize}
    \item Lmod spider walks entire module tree.
    \item Can build a reverse map from paths to modules
    \item Can map program \& libraries to modules.
    \item /opt/apps/i15/mv2\_2\_1/phdf5/1.8.14/lib/libhdf5.so.9
      $\Rightarrow$ phdf5/1.8.14(intel/15.02:mvapich2/2.1)
    \item Also helps with function tracking.
    \item Tmod Sites can still use Lmod to build the reverse map.
  \end{itemize}
\end{frame}

% page 9
\begin{frame}{Protecting XALT: Python to C++}
    \begin{itemize}
      \item XALT 1 used python scripts
      \item It was difficult to protect Python from users in everytime
      \item Solution: LD\_LIBRARY\_PATH="@ld\_lib\_path@"
        PATH=/usr/bin:/bin C++-exec ...
      \item Everything that depends on PATH must be hard coded
    \end{itemize}
\end{frame}

% page 10
\begin{frame}{Using XALT Data}
  \begin{itemize}
    \item Targetted Outreach: Who will be affected
    \item Largemem Queue Overuse
    \item XALT and TACC-Stats
  \end{itemize}
\end{frame}

% page 11
\begin{frame}{Tracking Non-mpi jobs (I)}
  \begin{itemize}
    \item Originally we tracked only MPI Jobs
    \item By hijacking mpirun etc.
    \item Now we can use ELF binary format to track jobs
  \end{itemize}
\end{frame}

% page 12
\begin{frame}[fragile]
    \frametitle{ELF Binary Format Trick}
 {\small
    \begin{semiverbatim}
void myinit(int argc, char **argv)
\{
  /* ... */
\}
void myfini()
\{
  /* ... */
\}
  static __attribute__((section(".init_array")))
       typeof(myinit) *__init = myinit;
  static __attribute__((section(".fini_array")))
       typeof(myfini) *__fini = myfini;
    \end{semiverbatim}
}
\end{frame}

% page 13
\begin{frame}{Using the ELF Binary Format Trick}
  \begin{itemize}
    \item This C code is compiled and linked in through the hijacked linker
    \item It can also be used with \texttt{LD\_PRELOAD}
    \item We are using both...
  \end{itemize}
\end{frame}

% page 14
\begin{frame}{Challenges (I)}
  \begin{itemize}
    \item Only want to track some executables on compute nodes
    \item Do not want to get overwhelmed by the data. 
  \end{itemize}
\end{frame}

% page 15
\begin{frame}{Answers}
  \begin{itemize}
    \item XALT Tracking only when told to
    \item Compute node only by host name filtering
    \item Executable Filter based on Path
    \item Protection against closing stderr before fini.
    \item Site configurable!
  \end{itemize}
\end{frame}

% page 16
\begin{frame}{Path Filtering}
  \begin{itemize}
    \item Uses FLEX to compile in patterns
    \item Use regex expression to control what to keep and ignore.
    \item Three files containing regex patterns, converted to code.
    \item Accept List Tests: Track /usr/bin/ddt, /bin/tar
    \item Ignore List Tests: /usr/bin, /bin, /sbin, ...
  \end{itemize}
\end{frame}

% page 17
\begin{frame}[fragile]
    \frametitle{TACC\_config.py}
 {\small
    \begin{semiverbatim}
hostname_patterns = [
  ['KEEP', '^c[0-9][0-9][0-9]-[0-9][0-9][0-9]\..*']
  ]
path_patterns = [
    ['PKGS',  r'.*/python[0-9][^/][^/]*'],
    ['PKGS',  r'.*/R'],
    ['KEEP',  r'^/usr/bin/ddt'],
    ['SKIP',  r'^/usr/.*'],
    ['SKIP',  r'^/bin/.*'],
  ]
env_patterns = [
    [ 'SKIP', r'^MKLROOT=.*' ],
    [ 'SKIP', r'^MKL_DIR=.*' ],
    [ 'KEEP', r'^I_MPI_INFO_NUMA_NODE_NUM=.*'],
  ]
    \end{semiverbatim}
}
\end{frame}

% page 18
\begin{frame}{Protecting XALT (I): UTF8 Characters}
  \begin{itemize}
    \item Linux supports UTF8 Characters in file names, env. vars.
    \item Python supports UTF8 if you know what you are doing.
    \item Switch XALT to use prepared statements
    \item Where \texttt{query="INSERT INTO table VALUE(?,?)"}
    \item This prevent SQL injection: ``johnny drop tables;''
    \item Also supports UTF8 characters.
  \end{itemize}
\end{frame}

% page 19
\begin{frame}{Protecting XALT (II): Python to C++}
    \begin{itemize}
      \item Difficult to protect Python from users in every case
      \item Solution: LD\_LIBRARY\_PATH="@ld\_lib\_path@"
        PATH=/usr/bin:/bin C++-exec ...
      \item Everything that depends on PATH must be hard coded
      \item basename $\Rightarrow$ /bin/basename
      \item Unique install for each operating system.
      \item Certain programs aren't in the same place: \texttt{basename}
    \end{itemize}
\end{frame}

% page 20
\begin{frame}{Speeding up XALT 2}
  \begin{itemize}
    \item Minimal impact on jobs ($>$ 0.09 secs)
    \item Non-mpi jobs only produce end record
    \item Filter by job\_id after transport.
    \item No more than 100 executions per job\_id.  (Changable)
    \item A launcher job 5 million executions in two hours.
  \end{itemize}
\end{frame}

% page 21
\begin{frame}{Tracking R packages}
  \begin{itemize}
    \item XALT 2 can now track R package usage
    \item James McComb \& Michael Scott from IU developed the R part
    \item They do this by intercepting the ``imports''
    \item Still in preliminary stages
    \item Plan to support Python later.
  \end{itemize}
\end{frame}

% page 22
\begin{frame}{New program: xalt\_extract\_record}
  \begin{itemize}
    \item This program reads the watermark.
    \item Find out who built this program on what machine
    \item Find out what modules where used.
  \end{itemize}
\end{frame}

% page 23
\begin{frame}[fragile]
    \frametitle{Example of xalt\_extract\_record output}
 {\tiny
    \begin{semiverbatim}
****************************************
XALT Watermark: hello
****************************************
Build_Epoch              1510257139.4624
Build_LMFILES            /opt/apps/modulefiles/intel/17.0.4.lua:...
Build_LOADEDMODULES      intel/17.0.4:impi/17.0.3:python/2.7.13:TACC:...
Build_OS                 Linux 3.10.0-514.26.2.el7.x86_64
Build_Syshost            stampede2
Build_UUID               586d5943-67eb-480b-a2fe-35e87a1f22c7
Build_User               mclay
Build_compiler           icc
Build_date               Thu Nov 09 13:52:19 2017
Build_host               c455-011.stampede2.tacc.utexas.edu
XALT_Version             1.7.7-devel
    \end{semiverbatim}
}
\end{frame}


% page 24
\begin{frame}{Tracking at Scale}
  \begin{itemize}
    \item TACC's Stampede2 has 1000's of nodes. (\# 15)
    \item TACC's Maverick has 132 nodes
    \item A single 1 hour Job generated 5 million records on Maverick
    \item It took 4.5 days to ingest!
    \item Cannot record everything so now what?
    \item Sampling
  \end{itemize}
\end{frame}


% page 25
\begin{frame}{Sampling: Three kinds of executables}
  \begin{itemize}
    \item Scalar: Non-mpi executables (Ntasks $==$ 1)
    \item PKGS: Track internal package usage: Python, R, Matlab
    \item MPI Programs (Ntasks $>$ 1)
  \end{itemize}
\end{frame}

% page 26
\begin{frame}{Changes in XALT behavior}
  \begin{itemize}
    \item Generate a start record before main()
    \item Generate an end record after main()
    \item Many MPI jobs run until time-out
    \item Can use accounting to set endtime
    \item Scalar now only produce an end record!
    \item Means that XALT won't know about Scalars that abort.
    \item An OpenMP program is a scalar program.
  \end{itemize}
\end{frame}

% page 27
\begin{frame}[fragile]
    \frametitle{Sampling of Scalar Prgms}
    \begin{itemize}
      \item TACC is using the following for Scalar programs:
 {\small
    \begin{semiverbatim}
     0 - 300 secs   0.01% chance
   300 - 600 secs   1%    chance
   600 - inf secs 100%    chance.
    \end{semiverbatim}
}
      \item Run time and probability are recorded.
      \item Can estimate total behavior from sampled data
      \item Can now sample usage of cp, mv, tar, perl, awk, ...
    \end{itemize}
\end{frame}

% page 28
\begin{frame}[fragile]
    \frametitle{Limitations of XALT}
 {\small
   \begin{itemize}
     \item XALT only knows the number of cores (Ntasks) really
     \item XALT can not tell the difference between
       \begin{semiverbatim}
    #SBATCH -N 10
    scalar_program
and
    #SBATCH -N 10
    ibrun -o 0 -n 1 scalar_program &
    ibrun -o 1 -n 1 scalar_program &
    ibrun -o 2 -n 1 scalar_program &
    ...
       \end{semiverbatim}
     \item num\_cores $=$ 1, num\_nodes $=$ 10 everywhere
     \item XALT can only reliably report Core-Hours not Node-Hours
     \item Cannot use XALT for accounting!
   \end{itemize}
}
\end{frame}

% page 29
\begin{frame}{Conclusion}
  \begin{itemize}
    \item Lmod:
      \begin{itemize}
        \item Source: github.com/TACC/lmod.git, lmod.sf.net
        \item Documentation: lmod.readthedocs.org
      \end{itemize}
    \item XALT:
      \begin{itemize}
        \item Source: github.com/xalt/xalt.git, xalt.sf.net
        \item Documentation: xalt.readthedocs.org
      \end{itemize}
  \end{itemize}
\end{frame}

%
\begin{frame}{License}
\footnotesize
\textcopyright The University of Texas at Austin, \the\year\\
\vspace{0.5 cm}
This work is licensed under the Creative Commons Attribution Non-Commercial 3.0 Unported License. To view a copy of this license, 
visit http://creativecommons.org/licenses/by-nc/3.0/\\
\vspace{0.5 cm}
When attributing this work, please use the following text: 
"Designing and Administering Large-scale Systems", Texas Advanced Computing Center, \the\year. Available under a Creative Commons Attribution Non-Commercial 3.0 Unported License.
\newline
\newline
\includegraphics[width=1cm]{./figures/cc.png}
\end{frame}

\end{document}
